\section{Conclusion and limitations}
I fine-tuned ResNet-50 and ViT-Base on HAM10000 on a laptop GPU, with a lesion-level split, class-weighted loss and augmentation, and evaluated both on held-out lesions with imbalance-aware metrics (Section~\ref{sec:results}).
With this short training budget the two architectures reach similar balanced accuracy (\resnetTestBacc{} and \vitTestBacc{}), and the difference is within test-set noise; ResNet-50 trains about 2.6 times faster per epoch on the laptop GPU.
Both remain weak on melanoma, the most clinically important class.

The study has clear limitations.
It is a single run with one seed and one split, with only a few epochs and no hyperparameter search, so the comparison may change under a larger budget; Vision Transformers lack some of the inductive biases of CNNs and depend more on pre-training scale~\cite{dosovitskiy2021image}, so their relative standing is sensitive to these choices.
The rarest classes have few test images (9 for dermatofibroma, 21 for vascular lesions), so their recall estimates are noisy.
PyTorch's MPS backend is not fully deterministic, so a rerun can give slightly different numbers.
Finally, HAM10000 was collected at two sites and does not represent all skin types, imaging devices or clinical settings~\cite{tschandl2018ham10000}; test performance on its own held-out lesions says nothing about clinical use.

Natural next steps are repeating the run over several seeds, training longer at higher resolution, and evaluating on an external dataset.
